\subsection{Jalon 1 - Du verrou au \textit{lock-free} : le Ring Buffer et le principe Single-Writer}

Le Jalon 0 a mis en évidence que le coût dominant du \textit{matching engine} ne réside pas dans la complexité algorithmique du chemin d'exécution, on a une médiane à 42\,ns pour \texttt{restingOrder} mais dans la queue de distribution, qui s'étale sur plusieurs ordres de grandeur sous l'effet d'événements externes au code métier. Parmi ces événements, les changements de contexte imposés par l'ordonnanceur lorsqu'un \textit{thread} se bloque sur un verrou constituent une source que l'architecture elle-même peut supprimer. C'est l'objet de ce jalon : remplacer les deux \texttt{ArrayBlockingQueue} du pipeline par une structure \textit{lock-free} inspirée du \textit{Disruptor} de LMAX, où aucun \textit{thread} n'est jamais suspendu par le noyau en attente d'une ressource partagée.

\subsubsection{Le compteur de séquence : la classe Sequence}

\textbf{Un compteur, deux garanties distinctes.} La coordination entre producteur et consommateur repose entièrement sur des compteurs monotones appelés \textit{séquences}. Une séquence n'est jamais un index physique dans un tableau : c'est une position logique dans un flux potentiellement infini, qui ne décroît jamais. Cette distinction, qui peut sembler cosmétique, est ce qui permet à deux \textit{threads} de comparer leur progression respective sans jamais confondre j'ai écrit 5\,000 éléments et je suis à la position 5\,000 du tableau, alors que ce tableau n'en contient que 1\,024.

\begin{lstlisting}[language=Java, caption={Compteur de séquence Single-Writer}]
public class Sequence {

    @SuppressWarnings("unused")
    private long padding1, padding2, padding3, padding4,
                 padding5, padding6, padding7;

    private volatile long value;

    @SuppressWarnings("unused")
    private long padding8, padding9, padding10, padding11,
                 padding12, padding13, padding14;

    public long get() {
        return value;
    }

    public void set(long value) {
        this.value = value;
    }

    public long incrementAndGet() {
        long next = value + 1;
        value = next;
        return next;
    }
}
\end{lstlisting}

Deux propriétés méritent d'être distinguées explicitement, car elles répondent à des problèmes différents.

La première est l'\textbf{atomicité}. La méthode \texttt{incrementAndGet()} n'est pas atomique : elle lit \texttt{value}, calcule \texttt{value + 1}, puis écrit le résultat en deux temps distincts. Sur une variable partagée par plusieurs producteurs, cette séquence serait vulnérable à une \textit{race condition} classique, deux \textit{threads} lisant la même valeur initiale, calculant chacun le même successeur, et l'un écrasant silencieusement le travail de l'autre. C'est précisément le problème que résout le \textit{Compare-And-Swap} (CAS) employé par \texttt{AtomicLong}. Ce mécanisme est ici inutile parce que le \textbf{Single-Writer Principle} garantit qu'un seul \textit{thread} écrit une instance donnée de \texttt{Sequence} : le producteur possède le curseur de réclamation, chaque consommateur possède sa propre séquence de progression. Aucune écriture concurrente n'est possible, donc aucune course à arbitrer. C'est le fondement architectural qui permet de s'affranchir du coût du CAS, dont l'état de l'art rappelle qu'il reste significatif (5\,700\,ms contre 300\,ms pour un simple \textit{thread} sans contrainte, cf. tableau~\ref{tab:comparaison_sync}).

La seconde propriété est la visibilité, garantie par le mot-clé \texttt{volatile}. Le Single-Writer Principle élimine les problèmes liés à l'asymétrie entre la vitesse de production et de consommation, mais ne dit rien de ce que voient les \textit{threads} lecteurs. Sans \texttt{volatile}, le compilateur JIT ou le processeur seraient libres de conserver la valeur écrite dans un registre ou dans le cache d'un cœur, sans jamais la republier en mémoire principale : le consommateur pourrait alors lire indéfiniment une valeur périmée, sans qu'aucune \textit{race condition} au sens strict ne se soit produite. \texttt{volatile} établit une relation \textit{happens-before} entre l'écriture du propriétaire et toute lecture ultérieure via \texttt{get()}, exactement comme le faisait le couple \texttt{put()}/\texttt{take()} d'une \texttt{BlockingQueue} au Jalon 0 mais sans le verrou qui l'accompagnait.

\textbf{Le \textit{padding} contre le False Sharing.} Les quatorze champs \texttt{long} inutilisés qui encadrent \texttt{value} ne servent à rien fonctionnellement : ils ne sont jamais lus ni écrits. Leur seul rôle est d'occuper de l'espace mémoire, pour garantir qu'aucune autre donnée ne partage la ligne de cache contenant \texttt{value}. Le dimensionnement de ce \textit{padding} obéit à un calcul précis.

Une instance de Java ne contrôle pas son adresse de départ, rien ne garantit que les variables de padding soit sur la même ligne de cache que la variable significative. Si on sous-dimensionne le padding en n'occupant artificiellment que 64 octets d'un côté (avant ou après la variable significative), peut-être que les variables de padding seront placées sur une autre ligne de cache et laisseront de l'espace sur la ligne de cache de la variable significative et donc de l'espace pour d'autres variables. Cela n'éliminerait donc pas le problème de \textit{false sharing} et rendrait le padding inutile. Dans le pire cas, \texttt{value} (8 octets) peut démarrer immédiatement après le début d'une ligne de cache, laissant jusqu'à 63 octets de cette même ligne et symétriquement en aval. Pour garantir que la ligne de cache contenant \texttt{value} n'appartienne qu'à cette instance, quel que soit son désalignement réel, il faut absorber cette marge des deux côtés. Comme \texttt{value} y participe lui-même pour 8 octets, le \textit{padding} strictement nécessaire est de $64 - 8 = 56$ octets de chaque côté, soit sept champs \texttt{long} de part et d'autre. C'est exactement le dimensionnement retenu historiquement par la classe \texttt{Sequence} du projet LMAX Disruptor.

La figure~\ref{fig:jol-sequence} présente la disposition mémoire obtenue avec JOL après application du \textit{padding}.

\begin{figure}[H]
    \centering
    \includegraphics[width=\textwidth]{SequenceJOLAfterPadding.png}
    \caption{Disposition mémoire de \texttt{Sequence} après \textit{padding}, obtenue avec Java Object Layout}
    \label{fig:jol-sequence}
\end{figure}

Trois éléments ressortent de cette mesure. D'abord, l'instance passe de 24 octets avant \textit{padding} à 136 octets après, soit $12$ octets d'en-tête, $4$ octets d'alignement, $56$ octets de \textit{padding} amont, $8$ octets utiles et $56$ octets de \textit{padding} aval. Une instance de 24 octets pouvait tenir à deux ou trois exemplaires dans une même ligne de cache : c'est précisément le scénario de \textit{False Sharing} que le Jalon 0 rendait possible entre le curseur de réclamation et le curseur de publication, tous deux instanciés consécutivement à la construction du \texttt{RingBuffer}.

Ensuite, le \textit{layout} confirme, sur une seconde classe après les observations sur la classe\texttt{Order}, le comportement de réordonnancement de la JVM déjà documenté au Jalon 0 : les champs n'apparaissent pas dans leur ordre de déclaration \texttt{padding14}, déclaré en dernier dans le code source, occupe le premier emplacement après l'en-tête. Ce réordonnancement est sans conséquence ici, tous les champs de \textit{padding} étant interchangeables par construction, mais il rappelle qu'un dimensionnement calculé au plus juste (par exemple trois champs avant et deux après, pour atteindre exactement 64 octets) reposerait sur une hypothèse d'ordre que la JVM ne garantit pas.

\subsubsection{Le tampon circulaire : la classe RingBuffer}

\textbf{Structure et indexation par masque binaire.} Le \texttt{RingBuffer} est un tableau d'objets de taille fixe, pré-alloué à la construction, dont la capacité doit être une puissance de deux. Cette contrainte n'est pas arbitraire : elle permet de substituer à l'opérateur modulo qui requiert une division entière, un simple ET logique avec un masque binaire.

\begin{lstlisting}[language=Java, caption={Structure du Ring Buffer et validation de capacité}]
public class RingBuffer<T> {

    private final Object[] buffer;
    private final int mask;
    private final Sequence cursor = new Sequence(-1);
    private final Sequence publishedSequence = new Sequence(-1);

    private final Sequence[] gatingSequence;
    private int countGatingSequence = 0;

    public RingBuffer(int capacity, int numberConsumers) {
        if (capacity <= 0)
            throw new IllegalArgumentException(
                "Need a capacity that is strictly positive");
        if (!((capacity & (capacity - 1)) == 0))
            throw new IllegalArgumentException(
                "Need a capacity that is a power of 2");

        this.buffer = new Object[capacity];
        this.mask = capacity - 1;
        this.gatingSequence = new Sequence[numberConsumers];
    }

    public T get(long sequence) {
        if (sequence < 0)
            throw new IllegalArgumentException(
                "Cannot read at index strictly negative");
        return (T) this.buffer[(int) (sequence & mask)];
    }
}
\end{lstlisting}

Le test \texttt{(capacity \& (capacity - 1)) == 0} est l'idiome standard de vérification d'une puissance de deux : une telle valeur ne possède qu'un seul bit à 1, et lui soustraire 1 inverse tous les bits, de sorte que le ET logique des deux vaut nécessairement zéro. Cette garde est indispensable : sans elle, le masquage ne produirait pas d'erreur visible mais un résultat silencieusement incorrect, projetant certaines séquences vers des emplacements erronés du tableau.

La conversion \texttt{(int) (sequence \& mask)} mérite un mot sur le typage. La séquence est un \texttt{long} et non un \texttt{int}, car elle représente une position dans un flux illimité : à un débit de 10\,000 opérations par seconde, un \texttt{int} déborderait après environ deux jours et demi de fonctionnement continu ; à un million d'opérations par seconde, le débit cible de la version finale, il déborderait en trente-cinq minutes. Un tel débordement corromprait immédiatement toute la logique de comparaison de séquences. Le masque, en revanche, est un \texttt{int}, car l'index physique qu'il produit est borné par nature à l'intervalle $[0, capacity - 1]$. La conversion explicite vers \texttt{int} n'intervient qu'après le masquage, à un moment où la valeur est mathématiquement garantie petite : elle ne perd donc aucune information, mais reste nécessaire pour satisfaire le compilateur, qui ne peut pas déduire statiquement cette borne.

\textbf{Deux curseurs, deux rôles distincts.} La coordination repose sur la séparation de deux compteurs que l'on pourrait confondre. Le \texttt{cursor} enregistre la dernière séquence réclamée en écriture par le producteur ; \texttt{publishedSequence} enregistre la dernière séquence effectivement \textit{publiée}, c'est-à-dire dont l'écriture est terminée et dont la donnée est prête à être lue.

Cette séparation résout un problème d'antériorité des opérations de lecture et d'écriture, non de performance. Le cycle d'écriture d'un producteur se décompose en trois appels distincts : \texttt{next()} réclame une position, \texttt{set()} y écrit la donnée, \texttt{publish()} signale que l'écriture est achevée. Si le consommateur se fiait au \texttt{cursor} pour déterminer jusqu'où lire, il pourrait observer une séquence réclamée mais dont \texttt{set()} n'a pas encore terminé l'écriture, et lire ainsi un emplacement vide ou contenant la donnée d'un tour précédent de l'anneau. C'est la raison pour laquelle \texttt{getCursor()}, seule méthode exposée aux consommateurs, retourne \texttt{publishedSequence} et jamais \texttt{cursor}.

\begin{lstlisting}[language=Java, caption={Publication et exposition de la séquence visible}]
public void publish(long sequence) {
    this.publishedSequence.set(sequence);
}

/**
 * Retourne la séquence de PUBLICATION la plus récente (visible pour un
 * consommateur), et non le curseur de réclamation interne.
 */
public long getCursor() {
    return publishedSequence.get();
}
\end{lstlisting}

\textbf{Le \textit{wrap point} et le \textit{back pressure} sans verrou.} Le tampon étant circulaire et borné, le producteur doit s'assurer qu'il n'écrase pas une donnée qu'un consommateur n'a pas encore lue. La méthode \texttt{next()} matérialise cette contrainte par le calcul d'un \textit{wrap point} : la séquence dont la position physique sera écrasée par la séquence candidate, soit $candidate - capacity$.

\begin{lstlisting}[language=Java, caption={Réclamation d'une séquence et contre-pression par busy-spin}]
public long next() {
    long candidateSequence = cursor.get() + 1;
    long wrapPoint = candidateSequence - this.capacity();

    while (wrapPoint > this.minGating()) {
        // busy-spin : attente active du consommateur le plus lent
    }

    cursor.set(candidateSequence);
    return candidateSequence;
}

private long minGating() {
    if (countGatingSequence == 0)
        return Long.MAX_VALUE;
    long minGating = this.gatingSequence[0].get();
    for (int i = 1; i < this.countGatingSequence; i++) {
        long minCandidate = this.gatingSequence[i].get();
        if (minGating > minCandidate)
            minGating = minCandidate;
    }
    return minGating;
}
\end{lstlisting}

Tant que le \textit{wrap point} dépasse la progression du consommateur, le producteur boucle activement sans jamais céder son cœur au système d'exploitation. C'est la différence fondamentale avec le \textit{back pressure} de l'\texttt{ArrayBlockingQueue} du Jalon 0 : là où \texttt{put()} suspendait le producteur via le noyau, avec le changement de contexte et le délai de réveil que cela implique, le \textit{busy-spin} le maintient en état \texttt{RUNNABLE}, prêt à reprendre dès l'instruction suivant la libération de l'emplacement. Le comportement observable est identique (le producteur ralentit quand le consommateur ne suit pas) mais le mécanisme sous-jacent ne sollicite jamais l'ordonnanceur.

\textbf{Généralisation multi-consommateurs.} La méthode \texttt{minGating()} retourne le minimum des progressions de tous les consommateurs enregistrés, et non la progression d'un consommateur unique. Cette généralisation dépasse le besoin du pipeline actuel, qui n'associe qu'un consommateur par tampon, mais elle est conforme à l'architecture du \textit{Disruptor} de LMAX : le producteur ne peut avancer qu'à la vitesse du consommateur le plus lent, faute de quoi il écraserait des données que ce dernier n'a pas encore traitées. Cette structure supporte donc nativement l'ajout ultérieur d'un second consommateur, un module de journalisation ou de diffusion de \textit{market data} en parallèle du \textit{matching}, par exemple, sans modification de \texttt{next()}.

Le choix d'injecter les séquences de \textit{gating} depuis l'extérieur, via \texttt{addGatingSequence()}, plutôt que de les créer à l'intérieur du \texttt{RingBuffer}, découle directement du Single-Writer Principle : c'est le consommateur qui possède sa séquence et l'écrit ; le \texttt{RingBuffer} ne fait que la lire.

\subsubsection{Intégration au pipeline : du blocage au \textit{busy-spin}}

\textbf{La boucle de consommation.} Le passage de l'\texttt{ArrayBlockingQueue} au \texttt{RingBuffer} transforme la nature même de la boucle de consommation. Avec une file bloquante, \texttt{take()} exprimait à la fois l'attente et la récupération de la donnée en un seul appel. Avec un tampon circulaire, le consommateur doit lui-même déterminer jusqu'où il a le droit de lire, en comparant sa propre progression à la séquence publiée.

\begin{lstlisting}[language=Java, caption={Boucle de consommation par busy-spin dans MatchingWorker}]
while (true) {
    while (this.inputRingBufferOrderCommand.getCursor()
            <= this.sequenceInputRingBuffer.get()) {
        if (Thread.interrupted()) {
            Thread.currentThread().interrupt();
            throw new InterruptedException(
                "Interrupted while spinning on inputRingBufferOrderCommand");
        }
    }

    long nextToRead = this.sequenceInputRingBuffer.get() + 1;
    inputCommand = inputRingBufferOrderCommand.get(nextToRead);
    this.sequenceInputRingBuffer.set(nextToRead);

    if (inputCommand.equals(OrderGenerator.SHUTDOWN_SIGNAL)) {
        long sequence = this.outputRingBufferTrade.next();
        outputRingBufferTrade.set(sequence, MatchingWorker.SHUTDOWN_SIGNAL);
        outputRingBufferTrade.publish(sequence);
        shutdownPropagated = true;
        break;
    }

    trades = this.engine.process(inputCommand);

    for (Trade trade : trades) {
        long sequence = this.outputRingBufferTrade.next();
        this.outputRingBufferTrade.set(sequence, trade);
        this.outputRingBufferTrade.publish(sequence);
    }
}
\end{lstlisting}

L'ordre des opérations dans le corps de la boucle mérite attention : la donnée est lue \textit{avant} que la séquence du consommateur ne soit avancée. Avancer la séquence en premier reviendrait à déclarer au producteur que l'emplacement est libre alors que la lecture n'a pas encore eu lieu, ouvrant la possibilité que le producteur l'écrase entre les deux instructions.

\textbf{La détection d'interruption, rendue explicite.} La vérification de \texttt{Thread.interrupted()} à l'intérieur de la boucle d'attente compense une perte induite par le \textit{busy-spin}. Un \textit{thread} bloqué sur \texttt{take()} était placé par le noyau en état \texttt{WAITING} et levait une \texttt{InterruptedException} dès l'appel à \texttt{interrupt()} : l'arrêt propre était garanti par la JVM elle-même. Un \textit{thread} en \textit{busy-spin} ne quitte jamais l'état \texttt{RUNNABLE} et ignore purement et simplement toute interruption qui ne serait pas explicitement testée. Sans cette vérification, le seul moyen d'arrêter le pipeline serait le signal applicatif, et une interruption externe, lors d'un arrêt forcé de la JVM, par exemple. Son absence laisserait les \textit{threads} tourner indéfiniment.

La construction en deux temps mérite d'être expliquée. \texttt{Thread.interrupted()} consomme le drapeau d'interruption lors de sa lecture. Il est donc immédiatement repositionné par \texttt{Thread.currentThread().interrupt()} avant que l'exception ne soit levée. Ainsi, l'appelant peut encore observer l'état interrompu.

\textbf{Le signal d'arrêt conservé dans le flux de données.} Deux stratégies étaient envisageables pour l'arrêt du pipeline : conserver la \textit{poison pill} transitant par le tampon comme une donnée ordinaire, ou introduire un signal hors flux sous la forme d'un \texttt{volatile boolean} partagé. La première a été retenue, pour une raison de méthode plutôt que de technique : introduire un second canal de communication en parallèle du \texttt{RingBuffer} ajouterait une variable à la comparaison entre le Jalon 0 et le Jalon 1, alors que l'objectif est précisément d'isoler l'effet du seul changement de mécanisme de transport. C'est le même raisonnement qui avait présidé au choix d'\texttt{ArrayBlockingQueue} plutôt que \texttt{LinkedBlockingQueue} pour la baseline.

\textbf{Une garantie d'arrêt dégradée.} La propagation du signal d'arrêt en cas d'interruption anormale, assurée dans le bloc \texttt{finally} de \texttt{MatchingWorker}, subit toutefois une dégradation par rapport à la baseline. La version \texttt{BlockingQueue} utilisait \texttt{offer()}, méthode non bloquante qui échouait proprement en retournant \texttt{false} si la file était pleine. Le \texttt{RingBuffer} n'expose aucun équivalent non bloquant : \texttt{next()} entre en \textit{busy-spin} si le tampon est saturé, et ce blocage ne peut être intercepté par aucun \texttt{catch}. Cette limitation est acceptée en l'état. Le tampon de sortie n'a pas de raison d'être saturé au moment d'un arrêt anormal en fonctionnement nominal mais elle constitue une régression de robustesse assumée.

\textbf{Assemblage.} L'assemblage du pipeline dans \texttt{Main} ne change que par la nature des structures de communication ; la topologie à trois \textit{threads} et la séquence de démarrage et de jonction restent identiques à celles du Jalon 0.

\begin{lstlisting}[language=Java, caption={Assemblage du pipeline avec deux Ring Buffers}]
RingBuffer<OrderCommand> commandQueue = new RingBuffer<>(RING_BUFFER_CAPACITY);
RingBuffer<Trade> tradeQueue = new RingBuffer<>(RING_BUFFER_CAPACITY);

MatchingEngine engine = new MatchingEngine();
String outputCsvPath = buildOutputCsvPath();

try (Writer csvWriter = new FileWriter(outputCsvPath)) {
    Thread generatorThread = new Thread(
            new OrderGenerator(commandQueue, COMMAND_COUNT, TARGET_RATE_PER_SECOND),
            "order-generator");
    Thread matchingThread = new Thread(
            new MatchingWorker(commandQueue, tradeQueue, engine),
            "matching-worker");
    Thread persistenceThread = new Thread(
            new PersistenceWorker(tradeQueue, csvWriter),
            "persistence-worker");

    generatorThread.start();
    matchingThread.start();
    persistenceThread.start();

    generatorThread.join();
    matchingThread.join();
    persistenceThread.join();
}
\end{lstlisting}

Chaque tampon possède son propre couple producteur-consommateur et sa propre séquence de \textit{gating}, enregistrée par le consommateur dans son constructeur. La capacité de 1\,024 éléments est conservée à l'identique du Jalon 0, précisément pour que cette dimension ne constitue pas une variable supplémentaire dans la comparaison.

\subsubsection{Le prix du \textit{lock-free} : consommation CPU et rigueur de test}

Deux contreparties de ce jalon méritent d'être énoncées explicitement, car elles ne sont pas des défauts d'implémentation mais le prix structurel du modèle retenu.

La première est la \textbf{consommation CPU continue}. Les deux \textit{threads} consommateurs ne dorment jamais : en l'absence de données à traiter, ils occupent intégralement un cœur à interroger la séquence publiée. Là où une \texttt{BlockingQueue} libérait le processeur pendant les périodes creuses, le \textit{busy-spin} paie une occupation permanente en échange de la suppression du délai de réveil. Ce compromis est délibéré et cohérent avec les stratégies employées en production dans l'industrie. Il suppose de disposer d'autant de cœurs que de \textit{threads} critiques comme formulé dans l'état de l'art.

La seconde est une \textbf{perte de rigueur dans les tests de concurrence}. Le test d'interruption du Jalon 0 s'appuyait sur \texttt{Thread.getState()} pour détecter de façon fiable qu'un \textit{thread} était effectivement bloqué avant de l'interrompre. Un \textit{thread} en \textit{busy-spin} restant en permanence \texttt{RUNNABLE}, cette détection n'est plus possible : le test se rabat sur un délai fixe, approximation temporelle intrinsèquement moins déterministe. Cette dégradation de l'outillage de test est une conséquence directe du modèle \textit{lock-free}, rarement mentionnée dans la littérature qui en vante les mérites de performance, et qu'il paraît honnête de documenter au même titre que ses bénéfices.

\subsubsection{Mesure comparative de la latence de transit cross-thread}

\textbf{Pourquoi \textit{HdrHistogram} plutôt que JMH.} La grandeur à mesurer ici n'est pas la durée d'exécution d'une méthode, mais un délai de \textit{transit} : il commence dans le \textit{thread} producteur, au moment où une donnée devrait être émise, et se termine dans le \textit{thread} consommateur, au moment où elle est effectivement lue. Aucune méthode unique ne l'encadre, ce qui place cette mesure hors du périmètre de JMH. Un benchmark spécifique a donc été construit, au prix de devoir réimplémenter manuellement la phase de \textit{warmup} que JMH fournissait nativement, et en s'appuyant sur \textit{HdrHistogram}, la bibliothèque de Gil Tene, pour l'enregistrement des latences. Ce choix avance l'introduction d'HdrHistogram, initialement prévue au jalon de métrologie, mais il présente un bénéfice immédiat : contrairement aux paliers fixes de JMH, HdrHistogram donne accès au percentile \texttt{P99.995} exact, celui-là même que fixent nos exigences non-fonctionnelles.

\textbf{Correction de l'omission coordonnée.} Un harnais naïf horodaterait chaque message au moment où le producteur parvient effectivement à le publier. Si le transport bloque le producteur (file pleine, verrou contesté) l'horodatage se décalerait d'autant et la latence mesurée resterait faible : le système paraîtrait rapide au moment précis où il est en difficulté. C'est le biais de \textit{coordonated omission} décrit par Tene, et il annulerait l'intérêt de la mesure. La correction consiste à fixer a priori une cadence d'émission et à horodater chaque message avec son instant d'émission \textit{théorique}, $t_0 + i \times \Delta t$, et non avec l'instant réel de publication. Si le producteur prend du retard, les instants théoriques continuent d'avancer sans lui : le retard accumulé se reporte intégralement sur la latence des messages suivants, au lieu d'être silencieusement absorbé.

\textbf{Résultats.} Les mesures suivantes ont été obtenues sur la machine de référence (10 cœurs, macOS ARM64, OpenJDK 26.0.1), sur un million de messages par configuration, après 200\,000 messages de chauffe non enregistrés.

\begin{table}[H]
    \centering
    \small
    \begin{tabular}{|l|l|r|r|r|r|r|r|}
        \hline
        Cadence & Transport & \textbf{Moy.} & \textbf{P50} & \textbf{P99} & \textbf{P99.9} & \textbf{P99.995} & \textbf{Max} \\
        \hline
        \multirow{2}{*}{100\,k/s}
          & \texttt{ArrayBlockingQueue} & 4\,922 & 4\,459 & 9\,215 & 33\,567 & 438\,783 & 730\,111 \\
          & \texttt{RingBuffer}         & 213   & 125   & 1\,000 & 14\,799 & 84\,223  & 575\,999 \\
        \hline
        \multirow{2}{*}{500\,k/s}
          & \texttt{ArrayBlockingQueue} & 2\,369 & 1\,750 & 8\,959 & 37\,759 & 181\,759 & 275\,455 \\
          & \texttt{RingBuffer}         & 57\,338 & 125  & 50\,751 & 12\,533\,759 & 14\,016\,511 & 14\,114\,815 \\
        \hline
        \multirow{2}{*}{1\,M/s}
          & \texttt{ArrayBlockingQueue} & 13\,433 & 1\,334 & 97\,279 & 2\,516\,991 & 3\,184\,639 & 3\,229\,695 \\
          & \texttt{RingBuffer}         & 413   & 125   & 8\,127 & 31\,967 & 156\,415 & 187\,903 \\
        \hline
    \end{tabular}
    \caption{Latence de transit cross-thread par percentile, en ns. Un million de messages par configuration.}
    \label{tab:jalon1-transport-latency}
\end{table}

\textbf{Agrégation sur cinq exécutions indépendantes.} Face à l'instabilité de la queue de distribution observée entre runs, cinq exécutions complètes du benchmark ont été menées, et la médiane du \texttt{P99.995} à travers ces cinq runs est retenue comme estimateur. Cette démarche applique au niveau des runs le même principe qui justifie l'usage du percentile plutôt que de la moyenne au niveau des échantillons : un estimateur robuste aux valeurs aberrantes isolées, ici des plateaux de perturbation externe qui contaminent un run sur cinq en moyenne sans se reproduire au même endroit.

\begin{table}[H]
    \centering
    \begin{tabular}{|l|l|r|r|r|r|r|}
        \hline
        Cadence & Transport & \multicolumn{5}{c|}{\texttt{P99.995} par run (\textmu s), triés} \\
        \hline
        \multirow{2}{*}{100\,k/s}
          & \texttt{ArrayBlockingQueue} & 226,6 & 264,7 & \textbf{271,9} & 438,8 & 5677,1 \\
          & \texttt{RingBuffer}         &  84,2 & 136,4 & \textbf{143,5} & 161,4 &  523,3 \\
        \hline
        \multirow{2}{*}{500\,k/s}
          & \texttt{ArrayBlockingQueue} &  59,9 &  64,1 & \textbf{79,8}  & 125,6 &  181,8 \\
          & \texttt{RingBuffer}         &  48,8 &  49,6 & \textbf{61,5}  & 143,0 & 14016,5 \\
        \hline
        \multirow{2}{*}{1\,M/s}
          & \texttt{ArrayBlockingQueue} &  57,0 &  71,4 & \textbf{206,8} & 524,3 & 3184,6 \\
          & \texttt{RingBuffer}         &  42,8 &  57,2 & \textbf{156,4} & 223,1 &  521,2 \\
        \hline
    \end{tabular}
    \caption{\texttt{P99.995} de latence de transit sur cinq runs indépendants, triés par valeur. La médiane (troisième valeur, en gras) est retenue comme estimateur robuste.}
    \label{tab:jalon1-transport-latency-runs}
\end{table}

\textbf{Lecture de l'agrégat.} Trois observations se dégagent de cette consolidation.

\begin{itemize}
    \item Premièrement, le classement entre transports, brouillé run par run, redevient net une fois agrégé : la médiane inter-runs du \textit{Ring Buffer} est inférieure à celle de l'\texttt{ArrayBlockingQueue} aux trois cadences 143,5 contre 271,9\,\textmu s à 100\,k/s (facteur 1,9), 61,5 contre 79,8\,\textmu s à 500\,k/s (facteur 1,3), 156,4 contre 206,8\,\textmu s à 1\,M/s (facteur 1,3). L'avantage du \textit{lock-free} sur la queue extrême, invisible ou même inversé sur un run isolé (cf. Run\,C2, où l'\texttt{ArrayBlockingQueue} atteint 524\,\textmu s à 1\,M/s tandis que le \textit{Ring Buffer} reste à 57\,\textmu s), se confirme donc statistiquement, mais avec un facteur de gain 1,3 à 1,9 nettement plus modeste que le facteur 36 observé sur la médiane intra-run. La suppression du verrou accélère massivement le cas courant ; son effet sur le pire cas, dominé par des perturbations exogènes communes aux deux transports, est réel mais contenu.
    \item Deuxièmement, \textbf{la contamination par un événement isolé touche indifféremment les deux transports et les trois cadences}, sans motif décelable : sur les trente séries de mesures (cinq runs, trois cadences, deux transports), trois présentent un \texttt{P99.995} supérieur à la milliseconde, \texttt{ABQ}/100\,k au Run\,B1, ainsi que \texttt{Ring}/500\,k et \texttt{ABQ}/1\,M, ces deux dernières contaminations provenant du même run (Run\,A) soit une configuration sur dix environ. Aucune récurrence sur un couple cadence/transport donné ne se dégage, ce qui achève d'écarter l'hypothèse initiale d'une fragilité propre au \textit{Ring Buffer} à charge intermédiaire, et confirme la nature exogène et aléatoire de la perturbation.
    \item Troisièmement, \textbf{l'exigence non-fonctionnelle de \texttt{P99.995} reste hors d'atteinte sur l'agrégat}, y compris pour le \textit{Ring Buffer} : 143,5\,\textmu s à 100\,k/s, 61,5\,\textmu s à 500\,k/s, le point le plus proche de la cible d'une cinquantaine de microsecondes, et 156,4\,\textmu s à 1\,M/s. Le mécanisme de transport, dont la médiane est désormais fixée au plancher de résolution de l'instrument (cf. quantification à 41,67\,ns), n'est plus le facteur limitant de cet écart : il est intégralement porté par les perturbations exogènes que ce jalon ne visait pas à traiter, et que le jalons suivants adresse sous un angle distinct.
\end{itemize}

Cette consolidation sur cinq runs confirme, avec la prudence statistique que la contamination du pire cas impose, l'hypothèse centrale du Jalon 1 : le passage au \textit{lock-free} réduit massivement et de façon parfaitement reproductible la latence typique de transit, et réduit plus modestement, mais de façon statistiquement significative, sa queue extrême. Le résidu de cette dernière relevant désormais de l'environnement d'exécution plutôt que du mécanisme de transport lui-même.